\section{A quantitative stability estimate}
\label{sec:bounded-stability}

If a $\theta$-plurisubharmonic function exceeds a bounded solution only by a small amount in $L^1$, the excess is also uniformly small. To prove this, we restrict the density of the solution to a strict superlevel set of the difference and construct a subsolution for the restricted density. Domination bounds the height of the superlevel set in terms of the norm of the restricted density, and H\"older's inequality converts its volume into a power of the $L^1$ error.

Throughout this section $\theta$ is semipositive and big, and $p>1$ is fixed.

\Needspace{20\baselineskip}
\begin{theorem}\label{thm:bounded-one-sided-stability}\index{Stability!bounded potentials}
Fix $L,C_0>0$. Let $u\in\PSH(X,\theta)$ satisfy
\[
 -C_0\leq u\leq0,\qquad
 \mathrm{MA}_\theta(u)=f\,\mathrm{d}V,\qquad
 0\leq f\in L^p(\mathrm{d}V),\quad\|f\|_p\leq L.
\]
For every
\[
 0<\eta<\frac{p-1}{np+p-1}
\]
there is a constant $C$, depending only on the fixed data and $\eta$, such that every $v\in\PSH(X,\theta)$ with $v\leq0$ satisfies
\[
 \sup_X(v-u)_+
       \leq C\left(\int_X(v-u)_+\,\mathrm{d}V\right)^\eta.
\]
\end{theorem}

The comparison potential $v$ need not be bounded.

\begin{proof}
Replacing $v$ by $\max(v,u)$ leaves both sides unchanged, so we may assume
$-C_0\leq u\leq v\leq0$. Put $I=\int_X(v-u)\,\mathrm{d}V$. If $I=0$, the two canonical potentials agree everywhere and there is nothing to prove, so we assume $I>0$.

\Needspace{5\baselineskip}
Choose $q\in(1,p)$, to be specified at the end. For $\delta>0$ put
\[
 E_\delta=\{u<v-\delta\},\qquad
 g_\delta=\mathbf1_{E_\delta}f,\qquad
 m_\delta=\|g_\delta\|_q.
\]
If $m_\delta>0$, apply \cref{lem:big-lp-subsolution} to $g_\delta/m_\delta$. It gives $h\in\PSH(X,\theta)$ and $b>0$ such that
\[
 -1\leq h\leq0,\qquad
 \mathrm{MA}_\theta(h)\geq b\,g_\delta/m_\delta\,\mathrm{d}V.
\]
The constant $b$ does not depend on $\delta,f,u,v$, since the auxiliary density has $L^q$ norm one.

First suppose $0<m_\delta<b/2$, and take
\[
 t=(2m_\delta/b)^{1/n}\in(0,1),\qquad
 w=(1-t)v+th-\delta-C_0t.
\]
Because $h-v\leq C_0$, we have $w\leq v-\delta$, and hence
$\{u<w\}\subset E_\delta$. Positivity of bounded mixed products gives
\[
 \mathrm{MA}_\theta(w)\geq t^n\mathrm{MA}_\theta(h)
       \geq2g_\delta\,\mathrm{d}V.
\]
On $\{u<w\}$ the last measure equals $2\mathrm{MA}_\theta(u)$. Bounded strict domination in the semipositive big background (\cref{thm:semipositive-bounded-domination}) therefore implies $u\geq w$. Consequently
\[
 v-u\leq v-w
      =\delta+C_0t+t(v-h)\leq\delta+(C_0+1)t.
\]
If $m_\delta=0$, the measure of $u$ vanishes on $E_\delta$, and the same domination theorem with coefficient zero, applied to $u$ and $v-\delta$, gives $v-u\leq\delta$. If $m_\delta\geq b/2$, the bound $v-u\leq C_0$ gives the following inequality after an increase of its constant. Thus in every case
\begin{equation}\label{eq:stability-small-set}
 \sup_X(v-u)\leq\delta+C_1m_\delta^{1/n}.
\end{equation}

H\"older's inequality and Chebyshev's inequality give
\[
 m_\delta\leq\|f\|_p\,\mathrm{d}V(E_\delta)^{1/q-1/p}
       \leq L(I/\delta)^{(p-q)/(pq)}.
\]
Set $a=(p-q)/(npq)>0$. For $0<I\leq1$, choose
$\delta=I^{a/(1+a)}$. Substituting in \eqref{eq:stability-small-set} yields
\[
 \sup_X(v-u)\leq C_2 I^{\eta_q},
 \qquad \eta_q=\frac{p-q}{npq+p-q}.
\]
As $q\downarrow1$, the exponent $\eta_q$ tends to $(p-1)/(np+p-1)$. For the prescribed $\eta$ we choose $q>1$ so close to $1$ that $\eta_q>\eta$; then $I^{\eta_q}\leq I^\eta$ on $(0,1]$. When $I>1$, the bound $v-u\leq C_0$ gives the same conclusion after increasing $C$. This proves the theorem.
\end{proof}

\Needspace{14\baselineskip}
\begin{corollary}\label{cor:bounded-stability-uniform-convergence}
Fix $C_0>0$. Let $u_j,u\in\PSH(X,\theta)$ satisfy $-C_0\leq u_j,u\leq0$ and
\[
 \mathrm{MA}_\theta(u_j)=f_j\,\mathrm{d}V,\qquad
 \mathrm{MA}_\theta(u)=f\,\mathrm{d}V,
 \qquad \sup_j\|f_j\|_p+\|f\|_p<+\infty.
\]
If $u_j\to u$ in $L^1$, then $u_j\to u$ uniformly. More precisely, for every exponent $\eta$ as in the theorem,
\[
 \|u_j-u\|_\infty\leq C\|u_j-u\|_{L^1(\mathrm{d}V)}^\eta.
\]
\end{corollary}
\begin{proof}
Apply the theorem first with solution $u$ and comparison potential $u_j$, and then with solution $u_j$ and comparison potential $u$. The constants are uniform because of the common range and the common density bound. The larger of the two suprema is the uniform norm, and each integral of a positive part is at most the full $L^1$ difference.
\end{proof}

The corollary assumes that the limit $u$ already solves an equation with an $L^p$ density. In an existence proof, this must be established before the corollary can be applied.
